\chapter{Comparative Statics and the Envelope Theorem}

Once an optimization problem is solved, the economically interesting
question is almost always the next one: how does the solution respond when
a parameter shifts? If the wage rises, how much labor does the firm shed?
If a tax is introduced, how fast does profit fall? \emph{Comparative
statics} answers the first question by implicit
differentiation of the first-order condition; the \emph{envelope theorem} answers the second: a small parameter change
moves the optimal \emph{value} only through its direct effect on the
objective, the induced re-optimization contributes nothing at first
order. Hotelling's lemma, the workhorse identity of producer theory, is a
direct instance. This chapter drills both tools for unconstrained
problems; the constrained version, where the derivative of the value
function is the Lagrange multiplier, follows in the chapters on the method
of Lagrange.

\section{Core Rules}

\begin{rulebox}{Parametric problem, optimizer, value function}
Let $f(x,\alpha)$ be a smooth objective with choice variable $x$ and
parameter $\alpha$, and suppose an interior maximum exists for each
$\alpha$. The \emph{optimizer} and the \emph{value function} are
\[
  x^*(\alpha)=\operatorname*{arg\,max}_x f(x,\alpha), \qquad
  V(\alpha)=\max_x f(x,\alpha)=f\bigl(x^*(\alpha),\alpha\bigr).
\]
In consumer theory $V$ is the indirect utility; in producer theory it is the profit
function. Throughout we abbreviate
$f_x=\pd{f}{x}$, $f_\alpha=\pd{f}{\alpha}$,
$f_{xx}=\dfrac{\partial^{2}f}{\partial x^{2}}$,
$f_{x\alpha}=\pdd{f}{x}{\alpha}$.
\end{rulebox}

\begin{rulebox}{Comparative statics: differentiate the FOC}
The optimizer solves $f_x\bigl(x^*(\alpha),\alpha\bigr)=0$ identically in
$\alpha$. If the second-order condition is strict
($f_{xx}(x^*,\alpha)\neq0$), the implicit function theorem makes
$x^*(\alpha)$ smooth, and differentiating the FOC gives
\[
  f_{xx}\cdot\dd{x^*}{\alpha}+f_{x\alpha}=0
  \quad\Longrightarrow\quad
  \dd{x^*}{\alpha}=-\frac{f_{x\alpha}(x^*,\alpha)}{f_{xx}(x^*,\alpha)}.
\]
At a maximum $f_{xx}<0$, so
$\sgn\Bigl(\dd{x^*}{\alpha}\Bigr)=\sgn\bigl(f_{x\alpha}\bigr)$: the
optimizer moves in the direction in which the marginal payoff $f_x$ rises
with $\alpha$. At a minimum ($f_{xx}>0$) the sign flips.
\end{rulebox}

\begin{rulebox}{Envelope theorem (unconstrained)}
For $V(\alpha)=\max_x f(x,\alpha)$ with interior optimizer $x^*(\alpha)$,
\[
  \dd{V}{\alpha}
  =\underbrace{f_x\bigl(x^*(\alpha),\alpha\bigr)}_{=0\text{ (FOC)}}
   \cdot\dd{x^*}{\alpha}
   +f_\alpha\bigl(x^*(\alpha),\alpha\bigr)
  =f_\alpha\bigl(x^*(\alpha),\alpha\bigr).
\]
Only the \emph{direct} effect of $\alpha$ on the objective matters; the
indirect effect through the re-optimized choice vanishes because $f_x=0$
at the optimum. The same identity holds at an interior minimum.
Geometrically, $V$ is the upper envelope of the
family of curves $\alpha\mapsto f(x,\alpha)$, one for each fixed $x$,
hence the name.
\end{rulebox}

\begin{rulebox}{Envelope workflow}
To find $\dd{V}{\alpha}$ without differentiating an explicit formula for
$V$:
\begin{enumerate}[leftmargin=1.6em,itemsep=0.3ex]
  \item Write the objective $f(x,\alpha)$.
  \item Differentiate with respect to the \emph{parameter only}:
        compute $f_\alpha(x,\alpha)$, treating $x$ as fixed.
  \item Evaluate at the optimum $x=x^*(\alpha)$ (explicitly known or
        symbolic).
\end{enumerate}
For $\dd{x^*}{\alpha}$ use the comparative-statics formula
$-f_{x\alpha}/f_{xx}$ instead.
\end{rulebox}

\begin{rulebox}{Economic readings}
\begin{itemize}[leftmargin=1.6em,itemsep=0.3ex]
  \item \textbf{Output price.} $\pi(p)=\max_q\,[\,pq-C(q)\,]$ gives
        $\pi'(p)=q^*(p)$: the marginal profit of a price increase equals
        current output.
  \item \textbf{Per-unit tax or cost.} If the objective contains $-tq$,
        then $\pd{V}{t}=-q^*$: profit falls at the rate of current
        output.
  \item \textbf{Wage.} For $\pi(w)=\max_L\,[\,p\,g(L)-wL\,]$,
        $\pd{V}{w}=-L^*$: the marginal loss equals current employment.
  \item \textbf{Prices in utility.} In quasi-linear consumer problems
        the derivative of indirect utility with respect to a price is
        $-(\text{quantity consumed})$.
\end{itemize}
\end{rulebox}

\begin{rulebox}{Hotelling's lemma}
Let a price-taking firm with production function $g$ solve
\[
  \pi(p,w)=\max_{L\ge0}\,\bigl[\,p\,g(L)-wL\,\bigr].
\]
Then, by the envelope theorem,
\[
  \pd{\pi}{p}=g\bigl(L^*(p,w)\bigr)=y^*(p,w), \qquad
  \pd{\pi}{w}=-L^*(p,w):
\]
differentiating the profit function with respect to the output price
recovers the \emph{supply function}, and with respect to an input price
recovers \emph{minus the factor demand}, no re-optimization needed.
\end{rulebox}

\begin{rulebox}{Curvature of the value function}
If $f(x,\alpha)$ is convex in $\alpha$ for every fixed $x$, then
$V(\alpha)=\max_x f(x,\alpha)$ is convex: a pointwise maximum of convex
functions is convex. In particular, profit functions are convex in
prices, since $p\mapsto pq-C(q)$ is affine for each fixed $q$. Combined
with Hotelling's lemma this yields $\pi''(p)=\dd{q^*}{p}\ge0$, the law
of supply.
\end{rulebox}

\section{Worked Example}

\begin{worked}{A firm maximises $\pi(L;w)=10\sqrt L-wL$. Find $L^{*}(w)$, the value function $\pi^{*}(w)$, and verify the envelope theorem}
\wstep{Step 1. Solve the problem for a \emph{fixed} parameter.}
\[
  \pd{\pi}{L}=\frac{10}{2\sqrt L}-w=\frac{5}{\sqrt L}-w=0
  \quad\Longrightarrow\quad
  \sqrt L=\frac5w
  \quad\Longrightarrow\quad
  \boxed{L^{*}(w)=\frac{25}{w^{2}}}
\]
Second order: $\pi_{LL}=-\tfrac52L^{-3/2}<0$, a maximum. \checkmark

\wstep{Step 2. Comparative statics: how the optimum responds.}
\[
  \dd{L^{*}}{w}=-\frac{50}{w^{3}}<0 .
\]
Labour demand falls when the wage rises --- as it must. Note this was
obtained by differentiating the \emph{solution}; the alternative, valid when
$L^{*}$ cannot be written down, is to differentiate the first-order
condition implicitly:
\[
  \frac{5}{\sqrt L}-w=0
  \;\Longrightarrow\;
  -\tfrac52L^{-3/2}\dd{L}{w}-1=0
  \;\Longrightarrow\;
  \dd{L}{w}=\frac{-1}{\tfrac52L^{-3/2}}=-\tfrac25L^{3/2} ,
\]
which at $L=25/w^{2}$ equals $-\tfrac25\cdot\frac{125}{w^{3}}=-\frac{50}{w^{3}}$.
\checkmark\ Same answer, and the second method is the one that generalises.
\wtrap{The sign of $\mathrm{d}L^{*}/\mathrm{d}w$ is the answer to
``does labour demand slope down?''. Implicit differentiation of the FOC
delivers it as $-\pi_{Lw}/\pi_{LL}$, where the denominator is negative by
the second-order condition --- so the sign of the response is the
\emph{sign of the cross partial}. That is the general comparative-statics
result, and this example is its simplest instance.}

\wstep{Step 3. Build the value function.}
Substitute the optimiser back into the objective:
\[
  \pi^{*}(w)=10\sqrt{\frac{25}{w^{2}}}-w\cdot\frac{25}{w^{2}}
  =10\cdot\frac5w-\frac{25}{w}
  =\frac{50-25}{w}
  =\boxed{\frac{25}{w}}
\]

\wstep{Step 4. Differentiate the value function directly.}
\[
  \dd{\pi^{*}}{w}=-\frac{25}{w^{2}} .
\]

\wstep{Step 5. Now the envelope theorem, which claims a shortcut.}
The theorem says: to differentiate the \emph{optimised} value with respect
to a parameter, differentiate the objective with respect to that parameter
and evaluate at the optimum --- \textbf{ignoring} how the optimum moves:
\[
  \dd{\pi^{*}}{w}=\left.\pd{\pi}{w}\right|_{L=L^{*}}
  =-L^{*}=-\frac{25}{w^{2}} . \checkmark
\]
One line, and it agrees with Step~4.

\wstep{Step 6. Why the moving optimum can be ignored.}
The full chain rule would read
\[
  \dd{\pi^{*}}{w}
  =\underbrace{\pd{\pi}{L}}_{=\,0\ \text{at }L^{*}}\dd{L^{*}}{w}
   +\pd{\pi}{w} .
\]
The first term vanishes \emph{because $L^{*}$ satisfies the first-order
condition}. The optimum does move, and its movement does change output ---
but at an optimum the objective is flat in $L$, so to first order that
movement contributes nothing.
\wtrap{The envelope theorem holds only \emph{at} the optimum, and only for a
parameter's \emph{direct} effect. Applying it at a non-optimal point, or
forgetting that $\partial\pi/\partial w$ means the derivative holding $L$
fixed, both give wrong answers. And it says nothing about second
derivatives --- $\mathrm{d}^{2}\pi^{*}/\mathrm{d}w^{2}$ genuinely requires
$\mathrm{d}L^{*}/\mathrm{d}w$.}

\wstep{Step 7. Read the result as economics.}
$\mathrm{d}\pi^{*}/\mathrm{d}w=-L^{*}$ is \textbf{Hotelling's lemma}: the
rate at which maximised profit falls when the wage rises equals the amount
of labour employed. A firm using $25$ workers loses about $25$ euros of
profit per euro of wage increase --- to first order, before it has time to
re-optimise. The whole family of ``derivative of the value function equals
the quantity'' results (Shephard's lemma, Roy's identity, the shadow price
$\lambda$ of Chapter~8) are this same theorem.

\wstep{Step 8. Curvature.}
$\pi^{*}(w)=25/w$ has $\mathrm{d}^{2}\pi^{*}/\mathrm{d}w^{2}=50/w^{3}>0$:
the value function is \textbf{convex} in the wage. This is general --- a
maximised value function is convex in the parameters it is linear in ---
and it has content: a firm facing a volatile wage with mean $\bar w$ does
better than one facing $\bar w$ for certain, because it can adjust. The
convexity \emph{is} the value of flexibility.
\end{worked}

\section{Practice Problems}

\begin{problems}
  \item For $f(x,\alpha)=-(x-\alpha)^2$, find $x^*(\alpha)$ and
        $V(\alpha)$, and verify the envelope theorem.
  \item Maximize $f(x,\alpha)=\alpha x-x^2$. Find $x^*(\alpha)$ and
        $\dd{x^*}{\alpha}$ (also via the comparative-statics formula),
        compute $V(\alpha)$, and check $V'(\alpha)=x^*(\alpha)$.
  \item For $f(x,a,b)=ax-\dfrac{b}{2}x^2$ with $b>0$, find $x^*(a,b)$
        and both partial derivatives of $V(a,b)$; verify each against the
        envelope theorem.
  \item Sign questions. (a) At an interior maximum with
        $f_{x\alpha}<0$, what is the sign of $\dd{x^*}{\alpha}$?
        (b) At an interior \emph{minimum} with $f_{x\alpha}>0$?
  \item Maximize $\pi(q,p)=pq-\dfrac{q^3}{3}$ over $q>0$ (with $p>0$).
        Find $q^*(p)$ and $\dd{q^*}{p}$, compute $V(p)$, and verify
        $V'(p)=q^*(p)$ both directly and by the envelope theorem.
  \item For $\pi(q,p,c)=pq-cq^2$ with $c>0$: find $q^*$, then
        $\pd{V}{p}$ and $\pd{V}{c}$ by the envelope theorem, and verify
        $\pd{V}{p}=q^*$ by computing $V$ directly.
  \item Maximize $\pi(q,p)=p\sqrt{q}-q$ for $p>0$. Find $q^*(p)$ and
        $\dd{q^*}{p}$, compute $V(p)$, and verify the envelope theorem.
  \item Let $V(p)=\max_q\bigl(pq-\dfrac{q^4}{4}\bigr)$ for $p>0$. Find
        $V'(p)$ via the envelope theorem \emph{without} computing $V$
        explicitly.
  \item A firm with a per-unit cost parameter: $\pi(q,c)=10q-q^2-cq$
        with $0<c<10$. Find $q^*(c)$, the sign of $\dd{q^*}{c}$, and
        $\dd{V}{c}$ by the envelope theorem; verify the latter directly.
  \item Labor demand: $\pi(L,w)=6\sqrt{L}-wL$ with $w>0$. Find
        $L^*(w)$ and $\dd{L^*}{w}$; compute $V(w)$ directly and check it
        against the envelope theorem.
  \item Demand response: a quasi-linear consumer solves
        $\max_{c>0} u(c,p)$ with $u(c,p)=\ln c-pc$ and price $p>0$. Find
        $c^*(p)$ and $\dd{c^*}{p}$, compute the indirect utility $V(p)$,
        verify $V'(p)=-c^*(p)$, and interpret.
  \item For $f(x,\alpha)=\alpha\ln x-x$ with $\alpha>0$: find
        $x^*(\alpha)$ and $\dd{x^*}{\alpha}$, obtain $V'(\alpha)$ by the
        envelope theorem, and verify it directly.
  \item A worker chooses hours $h\ge0$ to maximize
        $u(h,w)=wh-\dfrac{h^2}{2}-\dfrac{h^3}{6}$ at wage $w>0$. Find
        $h^*(w)$ and the sign of $\dd{h^*}{w}$, and give $\dd{V}{w}$ in
        terms of $h^*$.
  \item \emph{Minimize} $f(x,\alpha)=\alpha\mathrm{e}^{-x}+x$ over
        $x\in\R$ for $\alpha>0$. Find $x^*(\alpha)$ and $V(\alpha)$, and
        check that the envelope theorem also holds at a minimum.
  \item An output tax: $\pi(q,p,t)=(p-t)q-q^2$ with $0<t<p$. Find
        $q^*(p,t)$ and $\pd{V}{t}$ by the envelope theorem; verify
        directly and interpret economically.
  \item Hotelling's lemma: a price-taking firm produces $y=\sqrt{L}$
        and solves $\pi(p,w)=\max_{L\ge0}\,[\,p\sqrt{L}-wL\,]$. Find
        $L^*(p,w)$, the supply $y^*(p,w)$, and $\pi(p,w)$, and verify
        both parts of Hotelling's lemma: $\pd{\pi}{p}=y^*$ and
        $\pd{\pi}{w}=-L^*$.
  \item Maximize $f(K,L;r,w)=4K^{1/4}L^{1/4}-rK-wL$ over $K,L>0$.
        (a) Without solving the first-order conditions, give
        $\pd{V}{r}$ and $\pd{V}{w}$ via the envelope theorem, in terms
        of $K^*$ and $L^*$. (b) Solve the FOCs for $K^*(r,w)$,
        $L^*(r,w)$, and $V(r,w)$, and confirm your answer to (a).
  \item Show: if $f(x,\alpha)$ is convex in $\alpha$ for every fixed
        $x$, then $V(\alpha)=\max_x f(x,\alpha)$ is convex. Deduce that
        the profit function $\pi(p)=\max_q\,[\,pq-C(q)\,]$ is convex in
        $p$, and interpret $\pi''(p)\ge0$ in the light of Hotelling's
        lemma.
\end{problems}

\section{Solutions}

\begin{solutions}
  \item FOC: $f_x=-2(x-\alpha)=0\Rightarrow x^*(\alpha)=\alpha$. SOC:
        $f_{xx}=-2<0$, a maximum. $V(\alpha)=f(\alpha,\alpha)=0$, so
        $V'(\alpha)=0$. Envelope: $f_\alpha=2(x-\alpha)$, evaluated at
        $x=x^*=\alpha$ this is $0=V'(\alpha)$. Verified. (The peak
        merely slides along; its height never changes.)

  \item FOC: $\alpha-2x=0\Rightarrow x^*=\alpha/2$. SOC: $f_{xx}=-2<0$.
        Directly $\dd{x^*}{\alpha}=\dfrac12$; via the formula, with
        $f_{x\alpha}=1$:
        $\dd{x^*}{\alpha}=-\dfrac{f_{x\alpha}}{f_{xx}}
        =-\dfrac{1}{-2}=\dfrac12$. Consistent.
        $V(\alpha)=\alpha\cdot\dfrac{\alpha}{2}-\dfrac{\alpha^2}{4}
        =\dfrac{\alpha^2}{4}$, so $V'(\alpha)=\dfrac{\alpha}{2}=x^*$.
        Envelope: $f_\alpha=x\big|_{x^*}=\alpha/2$. Match.

  \item FOC: $a-bx=0\Rightarrow x^*=a/b$. SOC: $f_{xx}=-b<0$. Value:
        $V=a\cdot\dfrac ab-\dfrac b2\cdot\dfrac{a^2}{b^2}
        =\dfrac{a^2}{b}-\dfrac{a^2}{2b}=\dfrac{a^2}{2b}$. Then
        $\pd{V}{a}=\dfrac ab=x^*$, which is the envelope value
        $f_a=x\big|_{x^*}$; and $\pd{V}{b}=-\dfrac{a^2}{2b^2}$, which
        matches the envelope value
        $f_b=-\dfrac{x^2}{2}\Big|_{x^*}=-\dfrac{a^2}{2b^2}$. Both
        verified.

  \item (a) $\dd{x^*}{\alpha}=-f_{x\alpha}/f_{xx}$ with
        $f_{x\alpha}<0$ and $f_{xx}<0$: the quotient of two negatives is
        positive, and the leading minus makes
        $\dd{x^*}{\alpha}<0$. (b) At a minimum $f_{xx}>0$; with
        $f_{x\alpha}>0$, $\dd{x^*}{\alpha}=-(+)/(+)<0$. In both cases
        the optimizer falls as $\alpha$ rises.

  \item FOC: $p-q^2=0\Rightarrow q^*(p)=\sqrt{p}$. SOC:
        $\pi_{qq}=-2q=-2\sqrt{p}<0$ at $q^*$, a maximum.
        $\dd{q^*}{p}=\dfrac{1}{2\sqrt{p}}>0$; via the formula,
        $-\pi_{qp}/\pi_{qq}=-1/(-2\sqrt p)=\dfrac1{2\sqrt p}$.
        Consistent. Value:
        $V(p)=p\sqrt p-\dfrac{p^{3/2}}{3}=\dfrac23\,p^{3/2}$, so
        $V'(p)=p^{1/2}=q^*(p)$. Envelope:
        $\pi_p=q\big|_{q^*}=\sqrt p$. Match.

  \item FOC: $p-2cq=0\Rightarrow q^*=\dfrac{p}{2c}$. SOC: $-2c<0$.
        Envelope: $\pd{V}{p}=\pi_p\big|_{q^*}=q^*=\dfrac{p}{2c}$ and
        $\pd{V}{c}=\pi_c\big|_{q^*}=-(q^*)^2=-\dfrac{p^2}{4c^2}$.
        Direct check:
        $V=p\cdot\dfrac{p}{2c}-c\cdot\dfrac{p^2}{4c^2}
        =\dfrac{p^2}{2c}-\dfrac{p^2}{4c}=\dfrac{p^2}{4c}$, so
        $\pd{V}{p}=\dfrac{p}{2c}=q^*$ (and
        $\pd{V}{c}=-\dfrac{p^2}{4c^2}$). Verified.

  \item FOC: $\dfrac{p}{2\sqrt q}-1=0\Rightarrow\sqrt{q^*}=\dfrac p2
        \Rightarrow q^*=\dfrac{p^2}{4}$. SOC:
        $\pi_{qq}=-\dfrac{p}{4}q^{-3/2}<0$, a maximum.
        $\dd{q^*}{p}=\dfrac p2>0$. Value:
        $V(p)=p\cdot\dfrac p2-\dfrac{p^2}4=\dfrac{p^2}4$, so
        $V'(p)=\dfrac p2$. Envelope:
        $\pi_p=\sqrt q\big|_{q^*}=\dfrac p2$. Match.

  \item FOC: $p-q^3=0\Rightarrow q^*=p^{1/3}$; SOC: $-3q^2<0$ at
        $q^*>0$. Envelope:
        $V'(p)=\pi_p\big|_{q^*}=q^*=p^{1/3}$, with no need for $V$
        itself. (Cross-check: $V(p)=p\cdot p^{1/3}\cdot\ldots$,
        explicitly, $V=p^{4/3}-\dfrac14p^{4/3}=\dfrac34\,p^{4/3}$ and
        indeed $V'(p)=p^{1/3}$.)

  \item FOC: $10-2q-c=0\Rightarrow q^*(c)=\dfrac{10-c}{2}>0$. SOC:
        $-2<0$. Then $\dd{q^*}{c}=-\dfrac12<0$: a higher unit cost
        lowers output. Envelope:
        $\dd{V}{c}=\pi_c\big|_{q^*}=-q^*=-\dfrac{10-c}{2}$. Direct:
        $V=(10-c)q^*-(q^*)^2=2(q^*)^2-(q^*)^2=(q^*)^2
        =\dfrac{(10-c)^2}{4}$, so
        $\dd{V}{c}=\dfrac{2(10-c)\cdot(-1)}{4}=-\dfrac{10-c}{2}$.
        Verified.

  \item FOC: $\dfrac{3}{\sqrt L}-w=0\Rightarrow\sqrt L=\dfrac3w
        \Rightarrow L^*(w)=\dfrac{9}{w^2}$. SOC:
        $\pi_{LL}=-\dfrac32L^{-3/2}<0$. Then
        $\dd{L^*}{w}=-\dfrac{18}{w^3}<0$: labor demand falls in the
        wage. Direct value:
        $V(w)=6\cdot\dfrac3w-w\cdot\dfrac9{w^2}
        =\dfrac{18}w-\dfrac9w=\dfrac9w$, so $V'(w)=-\dfrac9{w^2}$.
        Envelope: $\pi_w=-L\big|_{L^*}=-\dfrac9{w^2}$. Match.

  \item FOC: $\dfrac1c-p=0\Rightarrow c^*(p)=\dfrac1p$. SOC:
        $u_{cc}=-\dfrac1{c^2}<0$. Then
        $\dd{c^*}{p}=-\dfrac1{p^2}<0$: demand slopes down. Indirect
        utility: $V(p)=\ln\dfrac1p-p\cdot\dfrac1p=-\ln p-1$, so
        $V'(p)=-\dfrac1p=-c^*(p)$. Envelope:
        $u_p=-c\big|_{c^*}=-\dfrac1p$. Verified. Interpretation: a
        marginal price increase lowers optimized utility at exactly the
        rate of current consumption: the re-optimization of $c$ has
        no first-order effect.

  \item FOC: $\dfrac{\alpha}{x}-1=0\Rightarrow x^*(\alpha)=\alpha$.
        SOC: $f_{xx}=-\dfrac{\alpha}{x^2}<0$ for $\alpha>0$. Then
        $\dd{x^*}{\alpha}=1$. Envelope:
        $V'(\alpha)=f_\alpha\big|_{x^*}=\ln x^*=\ln\alpha$. Direct:
        $V(\alpha)=\alpha\ln\alpha-\alpha$, so
        $V'(\alpha)=\ln\alpha+1-1=\ln\alpha$. Match.

  \item FOC: $w-h-\dfrac{h^2}{2}=0\Rightarrow h^2+2h-2w=0\Rightarrow
        h^*(w)=-1+\sqrt{1+2w}$ (the positive root). SOC:
        $u_{hh}=-1-h<0$ for $h\ge0$. Differentiating,
        $\dd{h^*}{w}=\dfrac{1}{\sqrt{1+2w}}>0$: hours rise with the
        wage. (Formula cross-check: $u_{hw}=1$, so
        $\dd{h^*}{w}=-\dfrac{1}{u_{hh}}=\dfrac{1}{1+h^*}
        =\dfrac{1}{\sqrt{1+2w}}$, using $1+h^*=\sqrt{1+2w}$.)
        Envelope: $\dd{V}{w}=u_w\big|_{h^*}=h^*=\sqrt{1+2w}-1$.

  \item FOC: $-\alpha\mathrm{e}^{-x}+1=0\Rightarrow
        \mathrm{e}^{-x}=\dfrac1\alpha\Rightarrow x^*(\alpha)=\ln\alpha$.
        SOC: $f_{xx}=\alpha\mathrm{e}^{-x}>0$, so this is a minimum.
        Value:
        $V(\alpha)=\alpha\cdot\dfrac1\alpha+\ln\alpha=1+\ln\alpha$, so
        $V'(\alpha)=\dfrac1\alpha$. Envelope:
        $f_\alpha=\mathrm{e}^{-x}\big|_{x^*}
        =\mathrm{e}^{-\ln\alpha}=\dfrac1\alpha$. Match, the FOC kills
        the indirect effect at a minimum exactly as at a maximum.

  \item FOC: $(p-t)-2q=0\Rightarrow q^*=\dfrac{p-t}{2}>0$. SOC:
        $-2<0$. Envelope:
        $\pd{V}{t}=\pi_t\big|_{q^*}=-q^*=-\dfrac{p-t}{2}$. Direct:
        $V=(p-t)q^*-(q^*)^2=(q^*)^2=\dfrac{(p-t)^2}{4}$, so
        $\pd{V}{t}=-\dfrac{p-t}{2}$. Verified. Interpretation: a
        marginal tax increase destroys profit at the rate of current
        output; the firm's downward output adjustment
        ($\dd{q^*}{t}=-\tfrac12$) affects profit only at second order,
        because the firm is already optimizing.

  \item FOC: $\dfrac{p}{2\sqrt L}-w=0\Rightarrow
        \sqrt{L^*}=\dfrac{p}{2w}\Rightarrow
        L^*(p,w)=\dfrac{p^2}{4w^2}$. SOC:
        $\pi_{LL}=-\dfrac p4L^{-3/2}<0$. Supply:
        $y^*=\sqrt{L^*}=\dfrac{p}{2w}$. Profit function:
        $\pi(p,w)=p\cdot\dfrac{p}{2w}-w\cdot\dfrac{p^2}{4w^2}
        =\dfrac{p^2}{2w}-\dfrac{p^2}{4w}=\dfrac{p^2}{4w}$. Hotelling:
        $\pd{\pi}{p}=\dfrac{2p}{4w}=\dfrac{p}{2w}=y^*$ and
        $\pd{\pi}{w}=-\dfrac{p^2}{4w^2}=-L^*$. Both parts verified.
        (Note also
        $\dfrac{\partial^{2}\pi}{\partial p^{2}}=\dfrac{1}{2w}
        =\pd{y^*}{p}>0$: the profit function is convex in $p$ and the
        supply curve slopes upward.)

  \item (a) The prices enter only through $-rK-wL$, so
        $f_r=-K$ and $f_w=-L$; by the envelope theorem
        $\pd{V}{r}=-K^*<0$ and $\pd{V}{w}=-L^*<0$, no solving
        needed. (b) FOCs: $f_K=K^{-3/4}L^{1/4}-r=0$ and
        $f_L=K^{1/4}L^{-3/4}-w=0$. Dividing them:
        $\dfrac{L}{K}=\dfrac{r}{w}\Rightarrow L=\dfrac rwK$;
        multiplying them:
        $K^{-1/2}L^{-1/2}=rw\Rightarrow KL=\dfrac{1}{(rw)^2}$.
        Substituting $L=rK/w$:
        $\dfrac rwK^2=\dfrac{1}{r^2w^2}\Rightarrow
        K^2=\dfrac{1}{r^3w}\Rightarrow K^*=r^{-3/2}w^{-1/2}$, hence
        $L^*=\dfrac rwK^*=r^{-1/2}w^{-3/2}$. (SOC:
        $f_{KK}=-\tfrac34K^{-7/4}L^{1/4}<0$ and
        $\det\Hess f=\tfrac9{16}K^{-3/2}L^{-3/2}
        -\tfrac1{16}K^{-3/2}L^{-3/2}
        =\tfrac12K^{-3/2}L^{-3/2}>0$: negative definite, a maximum.)
        Since $(K^*L^*)^{1/4}=\bigl((rw)^{-2}\bigr)^{1/4}=(rw)^{-1/2}$
        and $rK^*=wL^*=(rw)^{-1/2}$,
        \[
          V(r,w)=\frac{4}{\sqrt{rw}}-\frac{1}{\sqrt{rw}}
          -\frac{1}{\sqrt{rw}}=\frac{2}{\sqrt{rw}}.
        \]
        Check: $\pd{V}{w}=2\cdot\bigl(-\tfrac12\bigr)r^{-1/2}w^{-3/2}
        =-r^{-1/2}w^{-3/2}=-L^*$ and
        $\pd{V}{r}=-r^{-3/2}w^{-1/2}=-K^*$. Confirmed.

  \item Let $\alpha_\lambda=\lambda\alpha_0+(1-\lambda)\alpha_1$ with
        $\lambda\in[0,1]$, and let $\bar x$ maximize
        $f(\,\cdot\,,\alpha_\lambda)$. Then
        \[
          V(\alpha_\lambda)=f(\bar x,\alpha_\lambda)
          \le\lambda f(\bar x,\alpha_0)+(1-\lambda)f(\bar x,\alpha_1)
          \le\lambda V(\alpha_0)+(1-\lambda)V(\alpha_1),
        \]
        the first inequality by convexity of $f(\bar x,\cdot)$ in
        $\alpha$, the second because $V(\alpha_i)\ge f(\bar x,\alpha_i)$
        by definition of the maximum. Hence $V$ is convex. For the
        profit function: for each fixed $q$ the map $p\mapsto pq-C(q)$
        is affine, hence convex, so $\pi(p)$ is convex. Where $\pi$ is
        twice differentiable, Hotelling's lemma gives
        $\pi'(p)=q^*(p)$, so $\pi''(p)=\dd{q^*}{p}\ge0$: supply cannot
        slope downward, the law of supply.
\end{solutions}

\section{Additional Problems}

\begin{problems}
  \item For $f(x,\alpha)=6\alpha x-3x^2$, find $x^*(\alpha)$ and
        $V(\alpha)$, and verify the envelope theorem.
  \item Maximize $\pi(q,p)=pq-q^2-4$ (fixed cost $4$). Find $q^*(p)$
        and $\dd{q^*}{p}$, compute $V(p)$, and check $V'(p)=q^*(p)$.
  \item Labor demand: $\pi(L,w)=10\sqrt{L}-wL$ with $w>0$. Find
        $L^*(w)$ and $\dd{L^*}{w}$, compute $V(w)$, and verify the
        envelope theorem.
  \item For $f(x,\alpha)=\alpha\ln x-2x$ with $\alpha>0$: find
        $x^*(\alpha)$ and $\dd{x^*}{\alpha}$, obtain $V'(\alpha)$ via
        the envelope theorem, and verify it directly.
  \item For $\pi(q,c)=20q-(1+c)q^2$ with $c>0$: find $q^*(c)$ and
        $\pd{V}{c}$ via the envelope theorem; verify directly.
  \item A firm produces $y=2\sqrt{L}$ and solves
        $\pi(p,w)=\max_{L\ge0}\,[\,2p\sqrt{L}-wL\,]$. Find $L^*$, the
        supply $y^*$, and $\pi(p,w)$, and verify both parts of
        Hotelling's lemma.
  \item Maximize $f(K,L;r,w)=8\sqrt{K}+8\sqrt{L}-rK-wL$ over
        $K,L>0$. Find $K^*$, $L^*$, and $V(r,w)$, and verify
        $\pd{V}{r}=-K^*$ and $\pd{V}{w}=-L^*$.
  \item A worker with utility $u(h;w,t)=(w-t)h-\dfrac{h^2}{2}$, where
        $0<t<w$ is a wage tax, chooses hours $h\ge0$. Find $h^*(w,t)$,
        the sign of $\dd{h^*}{t}$, and $\pd{V}{t}$; verify directly and
        interpret.

  \item \textbf{Without re-solving.} The load a rectangular beam of width
        $B$ and height $H$ can carry is proportional to $BH^{2}$. The beam
        is cut from a circular log of diameter $d$, so $B^{2}+H^{2}=d^{2}$.
        Solve the maximisation once for general $d$, and give the maximal
        load $S^{*}(d)$.
  \item Using the previous answer, estimate the percentage increase in
        carrying capacity when the log's diameter grows from $30$\,cm to
        $31$\,cm --- \emph{without} re-solving the optimisation.
  \item A firm minimises cost $C=2K+3L$ subject to producing
        $\sqrt{KL}=q$. Find the value function $C^{*}(q)$ and use the
        envelope theorem to say by roughly how much minimum cost rises when
        the output target increases from $10$ to $10.5$.
  \item For $\max U=xy$ subject to $2x+5y=M$, the multiplier is
        $\lambda^{*}=M/20$. Estimate the gain in utility when income rises
        from $100$ to $104$, and compare with the exact change.
\end{problems}

\section{Additional Solutions}

\begin{solutions}
  \item FOC: $6\alpha-6x=0\Rightarrow x^*=\alpha$. SOC: $f_{xx}=-6<0$.
        $V(\alpha)=6\alpha^2-3\alpha^2=3\alpha^2$, so $V'(\alpha)=6\alpha$.
        Envelope: $f_\alpha=6x\big|_{x^*}=6\alpha$. Match.

  \item FOC: $p-2q=0\Rightarrow q^*=\dfrac p2$. SOC: $-2<0$.
        $\dd{q^*}{p}=\dfrac12$. Value:
        $V(p)=p\cdot\dfrac p2-\dfrac{p^2}4-4=\dfrac{p^2}4-4$, so
        $V'(p)=\dfrac p2=q^*$. Verified. (The fixed cost shifts $V$
        down but changes neither $q^*$ nor $V'$.)

  \item FOC: $\dfrac{5}{\sqrt L}-w=0\Rightarrow\sqrt L=\dfrac5w
        \Rightarrow L^*=\dfrac{25}{w^2}$. SOC:
        $\pi_{LL}=-\dfrac52L^{-3/2}<0$. Then
        $\dd{L^*}{w}=-\dfrac{50}{w^3}<0$. Value:
        $V(w)=10\cdot\dfrac5w-w\cdot\dfrac{25}{w^2}
        =\dfrac{50}w-\dfrac{25}w=\dfrac{25}w$, so
        $V'(w)=-\dfrac{25}{w^2}=-L^*$, exactly the envelope value
        $\pi_w\big|_{L^*}$. Verified.

  \item FOC: $\dfrac{\alpha}{x}-2=0\Rightarrow x^*=\dfrac{\alpha}{2}$.
        SOC: $-\dfrac{\alpha}{x^2}<0$. Then
        $\dd{x^*}{\alpha}=\dfrac12$. Envelope:
        $V'(\alpha)=\ln x\big|_{x^*}=\ln\dfrac{\alpha}{2}$. Direct:
        $V(\alpha)=\alpha\ln\dfrac{\alpha}{2}-\alpha$, so
        $V'(\alpha)=\ln\dfrac{\alpha}{2}
        +\alpha\cdot\dfrac{2}{\alpha}\cdot\dfrac12-1
        =\ln\dfrac{\alpha}{2}+1-1=\ln\dfrac{\alpha}{2}$. Match.

  \item FOC: $20-2(1+c)q=0\Rightarrow q^*=\dfrac{10}{1+c}$. SOC:
        $-2(1+c)<0$. Envelope: $\pi_c=-q^2$, so
        $\pd{V}{c}=-(q^*)^2=-\dfrac{100}{(1+c)^2}$. Direct:
        $V=20\cdot\dfrac{10}{1+c}-(1+c)\cdot\dfrac{100}{(1+c)^2}
        =\dfrac{200}{1+c}-\dfrac{100}{1+c}=\dfrac{100}{1+c}$, so
        $\pd{V}{c}=-\dfrac{100}{(1+c)^2}$. Verified.

  \item FOC: $\dfrac{p}{\sqrt L}-w=0\Rightarrow\sqrt{L^*}=\dfrac pw
        \Rightarrow L^*=\dfrac{p^2}{w^2}$. SOC:
        $-\dfrac p2L^{-3/2}<0$. Supply:
        $y^*=2\sqrt{L^*}=\dfrac{2p}{w}$. Profit function:
        $\pi(p,w)=2p\cdot\dfrac pw-w\cdot\dfrac{p^2}{w^2}
        =\dfrac{2p^2}{w}-\dfrac{p^2}{w}=\dfrac{p^2}{w}$. Hotelling:
        $\pd{\pi}{p}=\dfrac{2p}{w}=y^*$ and
        $\pd{\pi}{w}=-\dfrac{p^2}{w^2}=-L^*$. Both parts verified.

  \item The problem separates. FOCs:
        $\dfrac{4}{\sqrt K}-r=0\Rightarrow K^*=\dfrac{16}{r^2}$ and
        $\dfrac{4}{\sqrt L}-w=0\Rightarrow L^*=\dfrac{16}{w^2}$. SOC:
        the Hessian is diagonal with entries $-2K^{-3/2}<0$ and
        $-2L^{-3/2}<0$, negative definite. Value:
        $V=8\cdot\dfrac4r+8\cdot\dfrac4w-r\cdot\dfrac{16}{r^2}
        -w\cdot\dfrac{16}{w^2}
        =\dfrac{32}r+\dfrac{32}w-\dfrac{16}r-\dfrac{16}w
        =\dfrac{16}r+\dfrac{16}w$. Then
        $\pd{V}{r}=-\dfrac{16}{r^2}=-K^*$ and
        $\pd{V}{w}=-\dfrac{16}{w^2}=-L^*$, as the envelope theorem
        predicts. Verified.

  \item FOC: $(w-t)-h=0\Rightarrow h^*=w-t>0$. SOC: $u_{hh}=-1<0$.
        Then $\dd{h^*}{t}=-1<0$: the tax reduces hours worked.
        Envelope: $u_t=-h$, so $\pd{V}{t}=-h^*=-(w-t)$. Direct:
        $V=(w-t)^2-\dfrac{(w-t)^2}{2}=\dfrac{(w-t)^2}{2}$, so
        $\pd{V}{t}=-(w-t)$. Verified. Interpretation: at the margin the
        worker's welfare falls at the rate of hours currently worked;
        the induced cut in hours has no first-order welfare effect,
        pure envelope logic.

  \item Maximise $BH^{2}$ with $H^{2}=d^{2}-B^{2}$, i.e. maximise
        $B(d^{2}-B^{2})$. The derivative $d^{2}-3B^{2}$ vanishes at
        $B=d/\sqrt3$, giving $H^{2}=\tfrac23 d^{2}$ and
        \[
          S^{*}(d)=B H^{2}=\frac{d}{\sqrt3}\cdot\frac{2d^{2}}{3}
                  =\frac{2d^{3}}{3\sqrt3}.
        \]
        The optimal beam is always $H=\sqrt2\,B$, whatever the log.
  \item $S^{*}$ is proportional to $d^{3}$, so
        $\dfrac{\mathrm{d}S^{*}/S^{*}}{\mathrm{d}d/d}=3$: the elasticity of
        capacity with respect to diameter is exactly $3$. A rise from $30$
        to $31$ is $\tfrac{1}{30}=3.33\%$, so capacity rises by about
        $3\times3.33\%=10\%$. (Exactly:
        $(31/30)^{3}-1=10.34\%$.) No re-optimisation was needed --- only the
        envelope-style derivative of the value function.
  \item From $\sqrt{KL}=q$, $KL=q^{2}$. Minimising $2K+3L$ on that
        hyperbola gives $2=\mu L$, $3=\mu K$ (with $\mu$ the multiplier on
        $KL=q^{2}$), so $K/L=3/2$; substituting $KL=q^{2}$ gives
        $L=q\sqrt{2/3}$, $K=q\sqrt{3/2}$ and
        \[
          C^{*}(q)=2q\sqrt{\tfrac32}+3q\sqrt{\tfrac23}=2\sqrt6\,q
          \approx4.899\,q .
        \]
        By the envelope theorem $\mathrm{d}C^{*}/\mathrm{d}q=2\sqrt6$, so
        raising the target from $10$ to $10.5$ costs about
        $0.5\times4.899\approx2.45$. Here $C^{*}$ is linear in $q$, so the
        estimate is exact.
  \item The envelope theorem gives
        $\mathrm{d}V^{*}/\mathrm{d}M=\lambda^{*}=M/20$, which at $M=100$ is
        $5$. So $\Delta V^{*}\approx5\times4=20$. Exactly:
        $V^{*}(M)=M^{2}/40$, so
        $V^{*}(104)-V^{*}(100)=\tfrac{10816-10000}{40}=20.4$. The
        multiplier gives the first-order estimate; the $0.4$ gap is the
        curvature of the value function.
\end{solutions}
